The performance and reliability of an agentic workflow depend on the reasoning capabilities and specialized training of the underlying Large Language Models (LLMs). In my framework, rather than using a single general-purpose model, I adopted a multi-model strategy where different LLMs are assigned to specific nodes based on their strengths.

During the initial characterization phase, I evaluated five local models served through the Ollama backend. These models were tested across different task types as summarized in Table~\ref{tab:llm_comparison}. This benchmarking was necessary to identify the optimal ensemble for the final production deployment on Triton.

\begin{table}[h]
  \centering
  \caption{Comparative analysis of local LLMs for agentic workflow tasks.}
  \label{tab:llm_comparison}
  \begin{tabular}{|l|c|l|l|}
    \hline
    \textbf{Model} & \textbf{Size} & \textbf{Core Strengths}   & \textbf{Workflow Role}          \\
    \hline
    DeepSeek-R1    & 7B            & Reasoning, Debugging      & RAG Evaluation (LLM-as-a-Judge) \\
    Llama3-Groq    & 8B            & Structured Formats (JSON) & Workflow Routing                \\
    Qwen 2.5       & 7B            & Coding and Mathematics    & C Selection Logic               \\
    Mistral        & 7B            & Summarization, Multi-task & Search Summarization            \\
    GPT-OSS        & 20B           & Code Density, Context     & NNI Script Generation           \\
    \hline
  \end{tabular}
\end{table}

\subsubsection{Task-Specific Model Assignment}

The rationale for model assignment is driven by the specific computational requirements of each node:

For high-density code generation, although distilled reasoning models like \textbf{DeepSeek-R1} (7B) offer highly structured logic, empirical testing revealed that \textbf{GPT-OSS} (20B) consistently provides superior results for generating complex NNI trial scripts, largely because its higher parameter count translates directly to broader knowledge of Python libraries and far fewer syntax hallucinations. 

In addition, concerning inference latency versus reasoning tokens, \textbf{GPT-OSS} demonstrated noticeably lower wall-clock execution times compared to \textbf{DeepSeek-R1}. This speed advantage is primarily due to the complete absence of a ``thinking'' (Chain-of-Thought) token overhead; while R1 must painstakingly generate internal reasoning tokens before outputting functional code, the 20B model forcefully generates the final script directly with extremely high zero-shot accuracy. For structured output and routing, the routing nodes demand strict JSON compliance, a role for which \textbf{Llama3-Groq} (8B) was specifically retained due to its exceptionally high instruction-following performance in demanding low-latency scenarios. 

Finally, for specialized evaluation, \textbf{DeepSeek-R1} is exclusively reserved for the \textit{LLM-as-a-Judge} role within the detailed evaluation suite. Its unique ability to explicitly articulate underlying reasoning remains fundamentally important for transparently interpreting complex quality metrics (such as faithfulness and relevancy), scenarios where the inherent ``thinking'' overhead is considered a highly acceptable trade-off for unmatched explainability.

This tiered LLM architecture ensures that the agentic logic remains reliable while optimizing the hardware resources of the development server.
